%%------------------------------------------------------------------------------
%%Präambel
%%------------------------------------------------------------------------------
%%Klasse
%%------------------------------------------------------------------------------
\documentclass[a5paper,landscape,DIV=calc,11pt,notitlepage]{scrartcl}
%%------------------------------------------------------------------------------
%%Packages
%%------------------------------------------------------------------------------
\usepackage[headsepline=no, footsepline=no]{scrlayer-scrpage}
\usepackage[T1]{fontenc}
\usepackage{libertinus}
\usepackage[ngerman]{babel} 
\usepackage[right=24pt,left=24pt, top=24pt, bottom=24pt,twoside=false]{geometry}
\usepackage{tabularx}
\usepackage[svgnames,RGB,table]{xcolor}
\usepackage{microtype}
\usepackage[german]{layout}
\usepackage{tcolorbox}
\usepackage{hyperref}
\hypersetup{colorlinks, citecolor=, linkcolor=, urlcolor=blue,pdfauthor={Sascha Christmann},pdftitle={G\={o}j\={u}-Ry\={u} Karate-D\={o} Reusrath - Prüfungsordnung - Auszug zum 1. Dan},pdfsubject={G\={o}j\={u}-Ry\={u} Karate-D\={o}},pdfproducer={MiKTeX},pdfcreator={pdflatex}}
\usepackage{enumitem}
\usepackage{tabto}
%%------------------------------------------------------------------------------
%%Vereinbarungen & Makros
%%------------------------------------------------------------------------------
\clearpairofpagestyles
%%------------------------------------------------------------------------------
\definecolor{BKBELT}{rgb}{0.078,0.071,0.071}
%%------------------------------------------------------------------------------
\tcbset{width=\textwidth,height=\textheight,right=12pt,left=12pt,colback=white,fonttitle=\bfseries}
%%------------------------------------------------------------------------------
%%Ende Präambel
%%------------------------------------------------------------------------------
\begin{document}
%%------------------------------------------------------------------------------
%%Karte_Schwarz
%%------------------------------------------------------------------------------
\begin{tcolorbox}[colback=white,colframe=BKBELT,coltitle=white,title=1. Dan: Schwarz]
	\renewcommand{\labelenumi}{\alph{enumi}.)}
	\renewcommand{\labelenumii}{\arabic*.}
	{\scriptsize 
	\begin{tabularx}{\textwidth}{lX}
		\textit{Vorausgesetzte Techniken}:& Zuki-Waza, Uke-Waza, Geri-Waza und Dachi-Waza der bis hierhin erlernten Kihon-Ido, Kata, Kata-Bunkai und Yakusoku-Kumite
	\end{tabularx}\\}
	\begin{enumerate}[nosep]
		\item \textit{\textbf{Kihon-Ido}}:
		\begin{enumerate}[nosep]
			\item Harai-Otoshi-Uke/Gyaku-Zuki J\={o}dan/Mae-Geri Ch\={u}dan\tabto{335pt}(Zenkutsu-Dachi)
			\item Age-Uke/Ren-Zuki J\={o}dan-Ch\={u}dan\tabto{335pt}(Sanchin-Dachi)
			\item Mawashi-Empi-Uke/Uraken-Uchi/Harai-Otoshi-Uke/Gyaku-Zuki Ch\={u}dan\tabto{335pt}(Shiko-Dachi 45°)
		\item Shuto-Uke/Mae-Geri Ch\={u}dan\tabto{335pt}(Neko-Ashi-Dachi)
		\end{enumerate}
		\vspace{12pt}
		\item \textit{\textbf{Kata}}:\tabto{200pt}\textit{\textbf{Kata\,Bunkai}}:
		\begin{enumerate}[nosep]
			\item Sanseru\tabto{200pt}Sanseru
		\end{enumerate}
		\vspace{12pt}
		\item \textit{\textbf{Partnerformen (2 aus 3)}}:
		\begin{enumerate}[nosep]
			\item Yakusoku-Kumite:\tabto{200pt}3 Kumite-Ura\,\&\,2 Nage-Waza (frei auszuwählen)
			\item Shiai-Kumite:\tabto{200pt}Jiyu-Kumite: 1 Minute \frqq Freikampf\flqq~(sportlich)\\\tabto{200pt}\textit{entfällt für Sportler über 35 Jahre}
			\item Goshin-Jitsu-Kumite:\tabto{200pt}Abwehr und Konter gegen Packen, Würgen und Halten\\\textit{(Jissen-Jitsu, Selbstverteidigung)}\tabto{200pt} von vorn, von der Seite und von hinten
		\end{enumerate}
		\vspace{12pt}
		\item \textit{\textbf{Weitere Verbandsvorgaben}}:
		\begin{enumerate}[nosep]
			\item Vorlage eines Kampfrichterlehrgangsnachweis
			\item Besuch eines Dansha-Lehrgangs der Stilrichtung G\={o}j\={u}-Ry\={u} Karate-D\={o}
			\item Unsere Abteilung empfiehlt die Teilnahme an Lehrgängen auf Landes-, oder Bundesebene, um den Horizont des eigenen Karate beständig zu erweitern
			\item Weitere Details in der Prüfungsordnung
		\end{enumerate}
	\end{enumerate}
\end{tcolorbox}
%%------------------------------------------------------------------------------
\end{document}
%%------------------------------------------------------------------------------
%%EOF
%%------------------------------------------------------------------------------